\section{Theorem 5.2.1.1(3). Embedding a finite theta stage}

The global theta space is built from finite stages. Each stage keeps its
original topology inside the global space, as stated on
\href{https://github.com/mccorvie/riemann-masullo/blob/ac7e1cde68ff1fac01ecbef3c0e9f48e46bc1f06/paper/the-music-of-the-primes.pdf#page=230}{p.~230 of the paper}.
On p.~231 Masullo tries to recover the finite stage by projecting a global
test function. Two tests with the same global theta function can have
projections with different theta functions, so the result depends on the
representative. We instead add one finite place at a time and recover the
old theta function from two averages of the new one.

Let \(F\) be a number field and \(S\) a finite set of places containing the
archimedean places. Write \(B_S=\mathcal S(\mathbf A_{F,S})\) for the
Schwartz--Bruhat test space. The \(S\)-units \(G_S=\mathcal O_{F,S}^{\times}\)
are the elements of \(F^\times\) that are units at every place outside \(S\).
For \(f\in B_S\) and \(x\in\mathbf A_{F,S}^{\times}\), put
\[
  \Phi_{F,S}(f)(x)=\sum_{a\in G_S} f(a_S x),
\]
where \(a_S\) denotes the components of \(a\) at the places in \(S\).
Set \(Q_S=B_S/\ker\Phi_{F,S}\) and
\(Q_F=\mathcal S(\mathbf A_F)/\ker\Phi_F\), both with their quotient
topologies. These are the source and target of (5.2.1.3) on p.~230,
where Masullo writes both spaces as \(\mathrm H_{0,\mathrm{cont}}\).
The global theta operator \(\Phi_F\) sums over \(F^\times\).
We also write \(\Phi_{F,S}(z)\) for the theta function of a class \(z\in Q_S\),
since all representatives of \(z\) give the same values.

On tests, the insertion \(\pi_S\) extends \(f\) by multiplying by the
indicators \(1_{\mathcal O_w}\) at \(w\notin S\), where \(\mathcal O_w\) is
the ring of integers in \(F_w\). We prove that the induced map
\(\bar\pi_S\colon Q_S\to Q_F\) is a linear topological embedding.
We check its equivariance at the end.

\subsection{The projection on test functions}

The kernel identity \(\pi_S^{-1}(\ker\Phi_F)=\ker\Phi_{F,S}\) gives an
injective map of quotients. To get a continuous inverse on its image, the
paper uses the projection \(\Pi_S\) onto the inserted tests from
Lemma~2.1.4.3, p.~73. The claim on p.~231 is that \(\Pi_S\) preserves
\(\ker\Phi_F\), which would let it act on theta classes. The counterexample in
\href{https://github.com/mccorvie/riemann-masullo/blob/ac7e1cde68ff1fac01ecbef3c0e9f48e46bc1f06/ERRATA.md#L583}{erratum I123}
shows that this fails for every choice of the functional in that construction.

Take \(F=\mathbf Q\), \(S=\{\infty\}\), and \(e=1_{\widehat{\mathbf Z}}\).
Choose an even smooth function \(\phi\) supported within \(1/4\) of \(\pm1\)
with \(\phi(1)=\phi(-1)=1\). For \(f=\phi\otimes e\), the test
\(f-f(\,\cdot\,/2)\) has zero global theta function, since division by \(2\)
permutes \(\mathbf Q^\times\). The projection on the finite factor sends
\(1_{2\widehat{\mathbf Z}}\) to \(c e\) for some \(c\in\mathbf C\) and fixes
\(e\). The projected test is \(h\otimes e\), with \(h(y)=\phi(y)-c\phi(y/2)\).

Evaluate its global theta function at the idele with real component \(x\)
and all finite components \(1\). Only nonzero integers contribute because
\(\mathbf Q\cap\widehat{\mathbf Z}=\mathbf Z\). At \(x=1\), the first term of
\(h(ax)\) contributes at \(a=\pm1\) and the second at \(a=\pm2\), giving
\(2(1-c)\). At \(x=2\), only the second term contributes, at \(a=\pm1\),
giving \(-2c\). These two values cannot both vanish.

\subsection{Recovering the old theta function}

Let \(T=S\cup\{v\}\) for a finite place \(v\notin S\). Write
\(U_v=\mathcal O_v^\times\), let \(\mu\) be Haar probability on \(U_v\), and choose a
uniformizer \(\varpi\) with \(\operatorname{ord}_v(\varpi)=1\). For a class
\(z\in Q_T\) and an \(S\)-idele \(x\), take the difference of two averages
\begin{equation}
  D_v(z)(x)=\int_{U_v}
    \bigl(\Phi_{F,T}(z)(x,u)-\Phi_{F,T}(z)(x,\varpi^{-1}u)\bigr)
    \,d\mu(u).
  \label{5-2-1-1-3-theta-embedding:average}
\end{equation}
The formula uses only theta values, so it is well defined on \(Q_T\).
The two orbits are compact, and for each fixed test the theta summands have
a common summable bound on each orbit. This allows termwise integration.
(\href{https://github.com/mccorvie/riemann-masullo/blob/ac7e1cde68ff1fac01ecbef3c0e9f48e46bc1f06/RiemannMasullo/Ch05_ThetaZSpaces/Sec2.lean#L3506-L3530}{Summation and integration in Lean}.)

The insertion \(\bar\pi_{S,T}\colon Q_S\to Q_T\) sends \([f]\) to
\([f\otimes 1_{\mathcal O_v}]\). On this test, a term indexed by
\(a\in G_T\) has \(j=\operatorname{ord}_v(a)\). Its new factor in the first
evaluation is \(1_{j\geq 0}\), and in the second it is \(1_{j\geq 1}\),
because multiplication by \(\varpi^{-1}\) lowers the valuation by one.
These factors are independent of \(u\in U_v\). As \(\mu(U_v)=1\), the
coefficient of \(f(a_S x)\) in the difference is
\[
  1_{j\geq 0}-1_{j\geq 1}=1_{j=0}.
\]
An element of \(G_T\) with valuation zero at \(v\) is a unit outside \(S\),
so the surviving indices are exactly \(G_S\). This gives
\(D_v([f\otimes 1_{\mathcal O_v}])=\Phi_{F,S}(f)\).
(\href{https://github.com/mccorvie/riemann-masullo/blob/ac7e1cde68ff1fac01ecbef3c0e9f48e46bc1f06/RiemannMasullo/Ch05_ThetaZSpaces/Sec2.lean#L3924-L3944}{Indicator calculation in Lean}.)

\subsection{The left inverse on the larger space}

To use \eqref{5-2-1-1-3-theta-embedding:average} as a left inverse, we need
its output for every larger class to be a theta function from \(Q_S\). On a
product test \(f\otimes g\), with \(g\in\mathcal S(F_v)\), the coefficient of
\(f(a_S x)\) depends only on \(j=\operatorname{ord}_v(a)\). Write the
\(v\)-component of \(a\) as \(\varpi^j u_a\) with \(u_a\in U_v\).
The substitution \(u\mapsto u_a u\) preserves Haar measure, so the
coefficient is
\[
  w_j(g)=\int_{U_v}
    \bigl(g(\varpi^j u)-g(\varpi^{j-1}u)\bigr)\,d\mu(u).
\]
For sufficiently negative \(j\), both arguments lie outside the compact
support of \(g\). For sufficiently positive \(j\), both lie in a
neighborhood of zero where \(g\) is constant. Only finitely many
coefficients are therefore nonzero.

For each valuation \(j\) that occurs in \(G_T\), choose a representative
\(a_j\in G_T\). The elements of valuation \(j\) are exactly \(a_j b\) for
\(b\in G_S\), since \(G_S\) is the kernel of the valuation homomorphism.
Their contribution is
\[
  w_j(g)\sum_{b\in G_S}f(b_S a_{j,S}x)
      = w_j(g)\,\Phi_{F,S}(f)(a_{j,S}x).
\]
This is the theta function of the test \(w_j(g) f(a_{j,S}\,\cdot\,)\).
The full difference is a finite sum of such theta functions, taken over the
valuations that occur and have nonzero coefficient.
(\href{https://github.com/mccorvie/riemann-masullo/blob/ac7e1cde68ff1fac01ecbef3c0e9f48e46bc1f06/RiemannMasullo/Ch05_ThetaZSpaces/Sec2.lean#L3827-L3846}{The calculation for product tests in Lean}.)

Multiplication of tests identifies \(B_T\) with the algebraic tensor product
of \(B_S\) and \(\mathcal S(F_v)\), endowed with the inductive tensor topology
of Corollary~2.1.4.2, p.~72. Every test is a finite sum of product tests, so
the calculation gives a class in \(Q_S\) for every \(z\in Q_T\). Its theta
values \(D_v(z)\) determine the class uniquely. Since \(D_v\) is linear, so
is the resulting map \(r_v\colon Q_T\to Q_S\).

For fixed \(g\), the map \(f\mapsto r_v([f\otimes g])\) is a finite linear
combination of continuous translations followed by the quotient map. For
fixed \(f\), the map in \(g\) is linear and hence continuous because
\(\mathcal S(F_v)\) has the finest locally convex topology
(Proposition~2.1.2.8(1), p.~65). Separate continuity gives continuity on the
inductive tensor product \(B_T\), and the quotient topology gives continuity
of \(r_v\).
(\href{https://github.com/mccorvie/riemann-masullo/blob/ac7e1cde68ff1fac01ecbef3c0e9f48e46bc1f06/RiemannMasullo/Ch05_ThetaZSpaces/Sec2.lean#L3198-L3228}{Continuity on the quotient in Lean}.)

The indicator calculation gives \(r_v\bar\pi_{S,T}=\mathrm{id}_{Q_S}\).
Restricting \(r_v\) to the image gives a continuous inverse for the
insertion, so \(\bar\pi_{S,T}\) is a topological embedding.

\subsection{The topology of the global union}

Composing one-place embeddings gives all finite-stage embeddings. To recover
the stage topology inside the union, we extend its continuous seminorms.
By Proposition~2.1.4.1 every global test comes from a finite stage, so every
class in \(Q_F\) has a finite-stage representative. The compatible injections
thus identify \(Q_F\) with the union. A seminorm on \(Q_F\) is continuous
exactly when its pullback to the global test space is continuous. By the
final locally convex topology on tests, we can check this on each \(B_S\),
and by the quotient topology on \(Q_S\), this amounts to continuity on each
theta stage.
(\href{https://github.com/mccorvie/riemann-masullo/blob/ac7e1cde68ff1fac01ecbef3c0e9f48e46bc1f06/RiemannMasullo/Ch05_ThetaZSpaces/Sec2.lean#L2458-L2499}{Exhaustion and final topology in Lean}.)

Choose an increasing chain of finite stages starting at \(S\) and eventually
containing each finite set of places. Such a chain exists because the places
are countable. It gives the same final topology, since a seminorm continuous
on the chain restricts continuously to any finite stage by passing through a
chain stage containing it.

Start with a continuous seminorm \(p\) on \(Q_S\) and extend it along the
chain. At each step the old stage is a subspace \(M\) of the new stage.
Because \(p\) is continuous on \(M\), a suitably scaled finite maximum
\(P\) of continuous seminorms on the new stage has
\(\{P<1\}\cap M\subseteq\{p<1\}\). By homogeneity, \(p\leq P|_M\). Set
\[
  \widetilde p(y)=\inf_{m\in M}\bigl(p(m)+P(y-m)\bigr).
\]
Adding and scaling the candidates in the infimum gives a seminorm. Taking
\(m=0\) gives \(\widetilde p\leq P\), so it is continuous. For \(y\in M\) the triangle
inequality gives \(p(y)\leq p(m)+P(y-m)\) for all \(m\in M\), and taking
\(m=y\) gives equality in the infimum.
(\href{https://github.com/mccorvie/riemann-masullo/blob/ac7e1cde68ff1fac01ecbef3c0e9f48e46bc1f06/RiemannMasullo/ForMathlib/Analysis/LocallyConvex/WithSeminorms.lean#L19-L46}{Seminorm extension in Lean}.)

The extensions agree on overlaps and give a seminorm on the union, continuous
by finality and equal to \(p\) on \(Q_S\). Every basic neighborhood
\(\{p<1\}\) of zero in \(Q_S\) is therefore the restriction of a
neighborhood in \(Q_F\). Together with continuity of \(\bar\pi_S\), this
gives equality of the two topologies.
(\href{https://github.com/mccorvie/riemann-masullo/blob/ac7e1cde68ff1fac01ecbef3c0e9f48e46bc1f06/RiemannMasullo/Ch05_ThetaZSpaces/Sec2.lean#L2504-L2591}{Passage to the global embedding in Lean}.)

\paragraph{Scope and verification.}
Translation by an idele that is a unit
outside \(S\) fixes the outside integer indicators. Translating an inserted
test therefore agrees with translating its \(S\)-factor and then inserting.
The equality passes to theta classes.
(\href{https://github.com/mccorvie/riemann-masullo/blob/ac7e1cde68ff1fac01ecbef3c0e9f48e46bc1f06/RiemannMasullo/Ch05_ThetaZSpaces/Sec2.lean#L3174-L3183}{Equivariance on representatives in Lean}.)

Section~5.2.3, pp.~234--235, uses stage embeddings after passage to character
isotypes, which requires an argument on the character quotients. For the
ordinary theta spaces, Lean also proves that the image of \(Q_S\) in \(Q_F\)
is closed, as used in Lemma~5.3.1.5. It uses the embedding above and
completeness of \(Q_S\), obtained from the strict LF proof of part~(1).
(\href{https://github.com/mccorvie/riemann-masullo/blob/ac7e1cde68ff1fac01ecbef3c0e9f48e46bc1f06/RiemannMasullo/Ch05_ThetaZSpaces/Sec2.lean#L4210-L4221}{Closed-range argument in Lean}.)

The one-place repair comes from
\href{https://github.com/mccorvie/riemann-masullo/pull/4511}{PR~\#4511}.
The strict LF proof in
\href{https://github.com/mccorvie/riemann-masullo/pull/7176}{PR~\#7176}
reports an axiom check with no \texttt{sorryAx} in the
\href{https://github.com/mccorvie/riemann-masullo/blob/ac7e1cde68ff1fac01ecbef3c0e9f48e46bc1f06/RiemannMasullo/Ch05_ThetaZSpaces/Sec2.lean#L4128-L4132}{one-place left inverse}
or the
\href{https://github.com/mccorvie/riemann-masullo/blob/ac7e1cde68ff1fac01ecbef3c0e9f48e46bc1f06/RiemannMasullo/Ch05_ThetaZSpaces/Sec2.lean#L4205-L4208}{global embedding}.
